% Generated draft: requires mathematical and source review.
\documentclass[11pt]{article}
\usepackage[margin=25mm]{geometry}
\usepackage{fontspec}
\setmainfont{Latin Modern Roman}
\usepackage{amsmath,amssymb,mathtools,unicode-math}
\setmathfont{Latin Modern Math}
\usepackage{enumitem,needspace,xurl,hyperref}
\setlength{\parindent}{0pt}
\setlength{\parskip}{5pt}
\setlength{\emergencystretch}{3em}
\newcommand{\N}{\mathbb{N}}
\newcommand{\Z}{\mathbb{Z}}
\newcommand{\Q}{\mathbb{Q}}
\newcommand{\R}{\mathbb{R}}
\newcommand{\C}{\mathbb{C}}
\newcommand{\F}{\mathbb{F}}
\newcommand{\A}{\mathbb{A}}
\newcommand{\PP}{\mathbb{P}}
\newcommand{\E}{\mathbb{E}}
\newcommand{\eps}{\varepsilon}
\newcommand{\dd}{\,\mathrm{d}}
\begin{document}
\section*{UnsolvedMath ID: 20002022}
\textbf{Source problem code:} catalogue entry
\subsection{Definitions and mathematical statement}
Let $(M^6,g)$ be a closed oriented smooth Riemannian six-manifold. Normalize Branson's $Q$-curvature by $Q_6=120$ on the unit round $S^6$. Let $e_6(g)$ be the Euler scalar density relative to $\dd V_g$, normalized by $\int_M e_6(g)\,\dd V_g=\chi(M)$. A local conformal invariant $L$ has weight $-6$ and satisfies $L(e^{2\omega}g)=e^{-6\omega}L(g)$ pointwise. A divergence is $\operatorname{div}_gT(g)$ for a natural vector field and integrates to zero on $M$.

\textbf{Problem 4.} Can one explicitly express $Q_6$ as a constant times a Pfaffian, plus a local conformal invariant, plus a divergence, and in what sense is the expression unique? The original record immediately reports ``Robin Graham reports the answer to be YES.'' The trailing ``Alice Chang. General problems in conformal geometry'' begins subsequent material, rather than supplying another question here.

\textbf{Supplied formula and conventions.} The aid attributes to Graham, as recorded by Chang--Qing--Yang, the identity
\[
 Q_6=64\pi^3 e_6-\frac16\mathfrak J
                 +\frac1{10}\Delta^2R+\operatorname{div}T,
\]
where $\mathfrak J$ is a local conformal invariant. Fix $\Delta=\operatorname{div}\operatorname{grad}$ for the displayed current; then $\Delta^2R=\operatorname{div}(\nabla\Delta R)$, so
\[
 Q_6=64\pi^3e_6-\frac16\mathfrak J
             +\operatorname{div}\bigl(T+\tfrac1{10}\nabla\Delta R\bigr).
\]
The actual contracted formula for $\mathfrak J$ and the current $T$ is not supplied in this record, even though the aid describes it as explicitly available in the cited work. Thus this draft records the source's asserted decomposition and normalization, not an independently reconstructed fully explicit tensor formula.

For the standard Euler scalar $E_6$ with integral $384\pi^3\chi(M)$ and the AIM raw double-epsilon convention $\operatorname{Pff}_{\mathrm{raw}}=8E_6$, the same Euler summand is
\[
 64\pi^3e_6=\frac16E_6=\frac1{48}\operatorname{Pff}_{\mathrm{raw}}.
\]
Changing the Pfaffian convention without changing its coefficient would give a different formula.

\textbf{Uniqueness qualification from the aid.} Interpret the decomposition as a universal natural identity for all metrics in dimension six, with the Euler normalization fixed. The aid asserts uniqueness of the Euler coefficient and of the local-conformal-invariant \emph{class modulo conformal divergences}. More precisely, if $\mathcal I_6$ is the space of weight-$-6$ local conformal invariants and $\mathcal D_6$ the space of natural divergences, the determined class lies in $\mathcal I_6/(\mathcal I_6\cap\mathcal D_6)$. A representative can change by a divergence which is itself conformally invariant, compensated in the current; the current can further change by any divergence-free natural vector field. No canonical pointwise $L$ or $T$ is specified without a basis/normal-form convention and a current gauge. Therefore the source's word ``uniquely'' is not interpreted as literal uniqueness of every summand.
\subsection{Short English statement}
The source reports Graham's affirmative answer to the six-dimensional $Q$-curvature decomposition problem. The supplied aid fixes the Euler coefficient in three Pfaffian conventions and clarifies that uniqueness concerns the Euler coefficient and an invariant class modulo conformal divergences. The explicit invariant/current formulas are not included, and the representatives are not canonical.
\subsection{Sources}
\begin{itemize}
\item Local UnsolvedMath registry, permanent ID 20002022, source code catalogue entry; canonical \texttt{aim-geometry-notes.json}, notes[359].
\item AIM \emph{Conformal structure in geometry, analysis, and physics}, Problem 4 and reported affirmative answer by Robin Graham: \url{https://aimath.org/WWN/confstruct/confstruct.pdf}.
\item The research summary attributes the formula to Graham as recorded explicitly by Chang--Qing--Yang; it supplies no title, year, theorem number, URL or full bibliographic entry for that work.
\item Registry research attempt 1, dated 2026-08-14: partial result, problem status at run solved, rigor ``proved'', novelty confidence low. The machine-generated, non-peer-reviewed aid requires independent verification before citation.
\end{itemize}
\end{document}
